Both classes replicate the effect, with strongly class-dependent size (Table~\ref{tab:split_study}). For wrinkle ridges the protocol choice alone nearly doubles the reported score: mAP50 0.146 under the random split against 0.077 under the spatial split, for an identical model, budget, and data pipeline. The decomposition row separates the two contributions. Compared on \emph{identical} validation tiles, the random-trained model outscores the spatial-trained one (0.117 vs.\ 0.077): the only difference between the two is that the random model's training set contains tiles overlapping the validation blocks, so this $+0.040$ is a training-side contamination benefit, not an ability difference. The remaining $+0.029$ (0.146 vs.\ 0.117, same model, different validation tiles) is the evaluation-side component: the random protocol's own validation tiles share half their pixels with training tiles. For impact craters the total effect at this budget is much smaller (0.122 vs.\ 0.106): the near-saturated crater channel offers little memorisable local detail that is not already ubiquitous, so proximity to training tiles buys less.

Two practical consequences follow. First, scores reported under random tile splits on this dataset are not comparable to spatially split scores --- and the discrepancy is largest precisely for the classes where detection genuinely works, i.e.\ where there is signal to memorise. Development-time scores quoted at much larger training budgets under random-split protocols (Sect.~\ref{sec:yolo_ensemble}) are therefore expected to overstate generalisation increasingly with budget, as continued training rewards memorisation of near-duplicate samples that a spatial protocol excludes by construction. Second, the honest absolute level matters for the comparison of Sect.~\ref{sec:comp_ridge}: the dedicated single-class ridge detector under the spatial protocol (mAP50 0.077) is fully consistent with the committed six-class model's shared-protocol showing (MCC 0.076), confirming that YOLO's weakness on ridges is a property of the box representation on curvilinear features, not an artefact of any particular training run.
